\documentclass[11pt,a4paper]{article}

% --- geometry ---
\usepackage[margin=1in]{geometry}

% --- font ---
\usepackage{mathptmx}

% --- math ---
\usepackage{amsmath}
\usepackage{amssymb}
\usepackage{amsthm}
\usepackage{braket}

% --- diagrams ---
\usepackage{tikz}
\usepackage{quantikz}

% --- formatting ---
\usepackage{enumitem}
\usepackage{hyperref}
\usepackage{xcolor}

% --- colors ---
\definecolor{circuitblue}{RGB}{40,80,160}

% --- hyperref setup ---
\hypersetup{
    colorlinks=true,
    linkcolor=circuitblue,
    urlcolor=circuitblue,
    citecolor=circuitblue
}

% --- theorem environments ---
\theoremstyle{definition}
\newtheorem{definition}{Definition}[section]
\newtheorem{example}{Example}[section]

\theoremstyle{remark}
\newtheorem*{intuition}{Intuition}
\newtheorem*{keypoint}{Key Point}

% --- page style ---
\pagestyle{plain}

% --- no paragraph indent, use spacing instead ---
\setlength{\parindent}{0pt}
\setlength{\parskip}{6pt}

\begin{document}

% --- title block ---
\begin{center}
    {\LARGE \textbf{Lecture Title}}\\[8pt]
    {\large Quantum Machine Learning}\\[12pt]
    \hrule
\end{center}

\vspace{12pt}
\tableofcontents
\newpage

% --- content starts here ---

\section{Section Title}

Your content here.

\begin{definition}
A definition goes here.
\end{definition}

\begin{intuition}
An intuitive explanation goes here.
\end{intuition}

\begin{keypoint}
A key takeaway goes here.
\end{keypoint}

\begin{example}
A worked example goes here.
\end{example}

\end{document}
